% GENERATED FILE. Edit canonical lessons, then run npm run generate:books.
% canonical-source-hash: fnv1a64:ede1eb3477e605b3
% canonical-lessons: FR-C89-tortue, FR-C89-nouveau, FR-C89-vieux, FR-C89-rapide, FR-C89-facile

\chapter{Tortoise, New, Old, Fast, Easy}
\label{ch:fr-tortoise-new-old-fast-easy}
\hlchaptermodality{89}

\begin{chapteropening}
\textbf{By the end of this chapter:} \emph{I can name a tortoise, new, old, fast, easy.}

Complete the last lesson of chapter 89: I can name a tortoise, new, old, fast, easy.
\end{chapteropening}

\section[la tortue]{la tortue — a tortoise}

\label{lesson:FR-C89-tortue}



\begin{quote}
Before the new one: say the French for a frog, then the French for a tiger.
\end{quote}

\subsection*{You'll want to know: la tortue}
\textbf{la tortue} — \textquotedblleft{}a tortoise\textquotedblright{}.

\begin{cousinweb}
\textbf{tortue} $\leftarrow$ Old Provençal \emph{tortuga} — borrowed. (EtymDB 2.1) The usual account takes it to Late Latin \emph{tartaruchus}, \textquotedblleft{}of the underworld\textquotedblright{}, a name for a creature thought demonic.
\end{cousinweb}

\begin{grammarlens}[title={What you've built}]
One more word for this chapter.
\end{grammarlens}

\subsection*{Your turn}
\begin{itemize}
\raggedright
  \item \emph{Say it:} \emph{la tortue}
  \item \emph{Say it:} \emph{la tortue}, once more
  \item \emph{Say it:} say \emph{la tortue}
  \item \emph{Recall:} say the French for a frog, then the French for a tiger, then say \emph{la tortue} again
\end{itemize}

\begin{tcolorbox}[breakable,colback=teal!4,colframe=teal!35!black,title={Before you move on}]
What does la tortue mean? (\textquotedblleft{}A tortoise\textquotedblright{}.) Say it once more.
\end{tcolorbox}
\section[nouveau]{nouveau — new}

\label{lesson:FR-C89-nouveau}



\begin{quote}
Before the new one: say the French for a tiger, then the French for a tortoise.
\end{quote}

\subsection*{You'll want to know: nouveau}
\textbf{nouveau} — \textquotedblleft{}new\textquotedblright{}.

\begin{cousinweb}
\textbf{nouveau} $\leftarrow$ Old French \emph{novel} $\leftarrow$ Latin \emph{novellus}, \textquotedblleft{}new, young\textquotedblright{}, from \emph{novus}, \textquotedblleft{}new\textquotedblright{} — inherited. (EtymDB 2.1)
\end{cousinweb}

\begin{grammarlens}[title={What you've built}]
One more word for this chapter.
\end{grammarlens}

\subsection*{Your turn}
\begin{itemize}
\raggedright
  \item \emph{Say it:} \emph{nouveau}
  \item \emph{Say it:} \emph{nouveau}, once more
  \item \emph{Say it:} say \emph{nouveau}
  \item \emph{Recall:} say the French for a tiger, then the French for a tortoise, then say \emph{nouveau} again
\end{itemize}

\begin{tcolorbox}[breakable,colback=teal!4,colframe=teal!35!black,title={Before you move on}]
What does nouveau mean? (\textquotedblleft{}New\textquotedblright{}.) Say it once more.
\end{tcolorbox}
\section[vieux]{vieux — old}

\label{lesson:FR-C89-vieux}



\begin{quote}
Before the new one: say the French for a tortoise, then the French for new.
\end{quote}

\subsection*{You'll want to know: vieux}
\textbf{vieux} — \textquotedblleft{}old\textquotedblright{}.

\begin{cousinweb}
\textbf{vieux} $\leftarrow$ Old French \emph{viel} $\leftarrow$ Vulgar Latin \emph{*veclus} $\leftarrow$ Latin \emph{vetulus}, \textquotedblleft{}oldish\textquotedblright{}, from \emph{vetus}, \textquotedblleft{}old\textquotedblright{} — inherited. (EtymDB 2.1)
\end{cousinweb}

\begin{grammarlens}[title={What you've built}]
One more word for this chapter.
\end{grammarlens}

\subsection*{Your turn}
\begin{itemize}
\raggedright
  \item \emph{Say it:} \emph{vieux}
  \item \emph{Say it:} \emph{vieux}, once more
  \item \emph{Say it:} say \emph{vieux}
  \item \emph{Recall:} say the French for a tortoise, then the French for new, then say \emph{vieux} again
\end{itemize}

\begin{tcolorbox}[breakable,colback=teal!4,colframe=teal!35!black,title={Before you move on}]
What does vieux mean? (\textquotedblleft{}Old\textquotedblright{}.) Say it once more.
\end{tcolorbox}
\section[rapide]{rapide — fast}

\label{lesson:FR-C89-rapide}



\begin{quote}
Before the new one: say the French for new, then the French for old.
\end{quote}

\subsection*{You'll want to know: rapide}
\textbf{rapide} — \textquotedblleft{}fast\textquotedblright{}.

\begin{cousinweb}
\textbf{rapide} $\leftarrow$ Latin \emph{rapidus}, \textquotedblleft{}swift, snatching\textquotedblright{}, from \emph{rapere}, \textquotedblleft{}to seize\textquotedblright{} — a learned borrowing. (EtymDB 2.1)
\end{cousinweb}

\begin{grammarlens}[title={What you've built}]
One more word for this chapter.
\end{grammarlens}

\subsection*{Your turn}
\begin{itemize}
\raggedright
  \item \emph{Say it:} \emph{rapide}
  \item \emph{Say it:} \emph{rapide}, once more
  \item \emph{Say it:} say \emph{rapide}
  \item \emph{Recall:} say the French for new, then the French for old, then say \emph{rapide} again
\end{itemize}

\begin{tcolorbox}[breakable,colback=teal!4,colframe=teal!35!black,title={Before you move on}]
What does rapide mean? (\textquotedblleft{}Fast\textquotedblright{}.) Say it once more.
\end{tcolorbox}
\section[facile]{facile — easy}

\label{lesson:FR-C89-facile}



\begin{quote}
Before the new one: say the French for old, then the French for fast.
\end{quote}

\subsection*{You'll want to know: facile}
\textbf{facile} — \textquotedblleft{}easy\textquotedblright{}.

\begin{cousinweb}
\textbf{facile} $\leftarrow$ Latin \emph{facilis}, \textquotedblleft{}easy to do\textquotedblright{}, from \emph{facere}, \textquotedblleft{}to do, to make\textquotedblright{} — a learned borrowing. (EtymDB 2.1)
\end{cousinweb}

\begin{grammarlens}[title={What you've built}]
One more word for this chapter.
\end{grammarlens}

\subsection*{Your turn}
\begin{itemize}
\raggedright
  \item \emph{Say it:} \emph{facile}
  \item \emph{Say it:} \emph{facile}, once more
  \item \emph{Say it:} say \emph{facile}
  \item \emph{Recall:} say the French for old, then the French for fast, then say \emph{facile} again
\end{itemize}

\begin{tcolorbox}[breakable,colback=teal!4,colframe=teal!35!black,title={Before you move on}]
What does facile mean? (\textquotedblleft{}Easy\textquotedblright{}.) Say it once more.
\end{tcolorbox}
